\subsection{Even-codegree parity closes the multiplicity-nine Gauss case}
\label{subsec:gauss-eight-general}
Let $R(8)$ be the Gauss-family tensor with basis $1,a,b,c,d,e$,
canonically identified from the full tensor as $\texttt{R6\_9\_8200}$.
Its original formal codegrees are
\[
D=4010+200\sqrt{401},\quad D'=4010-200\sqrt{401},\quad5,5,5,5.
\]
Its Casimir trace is $8040$, its formal-codegree trace-ideal multiple is
$4010$, and its Casimir polynomial reduces modulo two to
$X^2(X+1)^4$.

\begin{lemma}[Even dimension-codegree isolation]
Suppose a commutative self-dual fusion ring has an original formal codegree
$f$ with $f/2$ integral. Let $h$ be the multiplicity of zero in the reduction
modulo two of its Casimir characteristic polynomial. If the ring is
categorifiable, the dimension-character field $K$ of its signed
sphericalization satisfies $[K:\mathbb Q]\le h$.
\end{lemma}
\begin{proof}
Use the signed sphericalization algebra $A$, dimension character
$d=(1,d_1,\ldots,d_{r-1})$, and $\Delta=\sum_i d_i^2$ of the
dimension-field containment lemma. Central induction gives
$\Delta/f\in\mathcal O_K$ and $f\in K$. Hence
$\Delta/2=(f/2)(\Delta/f)$ is integral.
The rational algebra $A\otimes\mathbb Q$ is a product of totally real
number fields, one factor being $K$. Multiplication by the Casimir element
on this factor has characteristic polynomial
\[
P_K(X)=\operatorname{Norm}_{K/\mathbb Q}(X-\Delta).
\]
Its coefficients are integers, even when $\Delta$ generates a proper
subfield of $K$. Since $\Delta\in2\mathcal O_K$,
$P_K(X)\equiv X^{[K:\mathbb Q]}\pmod2$.
Every other field factor also has an integral characteristic polynomial.
The full signed Casimir polynomial is their product, and is congruent
modulo two to the original one. Therefore $[K:\mathbb Q]\le h$.
\end{proof}

For $R(8)$, the minimal polynomial of $D/2$ is
$X^2-4010X+10025$, so the lemma applies with $h=2$.
Induction also gives $\mathbb Q(D)=E=\mathbb Q(\sqrt{401})\subseteq K$,
and consequently $K=E$. This argument does not assume that the original
category has a pivotal structure.

The maximal order of $E$ is $\mathbb Z[\omega]$, where
$\omega=(1+\sqrt{401})/2$ and
$\omega^2-\omega-100=0$. Thus
$\mathcal O_E/2\mathcal O_E\cong\mathbb F_2\times\mathbb F_2$.
Direct exhaustive evaluation of the $32$ possible nonunit coordinate
vectors gives exactly one unital character $R(8)/2R(8)\to\mathbb F_2$:
\[
(1,a,b,c,d,e)\longmapsto(1,0,0,1,0,0).
\]
The signed coefficient parity is the same as the original parity.
Composing the signed dimension character with either projection of
$\mathcal O_E/2\mathcal O_E$ therefore gives this same character. It follows
that $d_a,d_b,d_d,d_e\in2\mathcal O_E$ and
$d_c\in1+2\mathcal O_E$. Therefore
\[
\Delta=1+d_a^2+d_b^2+d_c^2+d_d^2+d_e^2
       \equiv2\pmod{4\mathcal O_E},
\qquad \operatorname{Tr}_{E/\mathbb Q}(\Delta)\equiv0\pmod4.
\]
The trace-ideal bound gives the separate divisibility by $4010$.
Consequently the positive trace of $\Delta$ is a multiple of
$\operatorname{lcm}(4010,4)=8020$. Since it is bounded by $8040$, it is
exactly $8020$. Hence the signed Casimir trace is at least $8020$ and its
deficit from the original trace is at most $20$.

The minimum possible loss from decreasing any signed coefficient magnitude
is $28$: the minimum over unordered nonunit triples with coefficient $a$
and orbit size $w$ is $w(a^2-(a-2)^2)=28$. Since $20<28$, every signed
coefficient has its original magnitude. The next lemma makes the final
sign check very small and imposes no Frobenius--Schur indicator condition.

\begin{lemma}[Repeated-index gauge under saturation]
Suppose $R$ is a commutative self-dual fusion ring and $T$ a commutative
self-dual associative signed table with $|t_{ijk}|=N_{ijk}$.
Suppose $N_{iik}>0$ for every nonunit $i,k$. Then a basis sign gauge makes
all repeated-index nonunit coefficients of $T$ positive.
\end{lemma}
\begin{proof}
Associativity and the coefficient-of-unit inner product give, for any
nonunit $i,j$,
\[
\sum_k t_{iik}t_{jjk}=\sum_k t_{ijk}^2
=\sum_k N_{ijk}^2=\sum_k N_{iik}N_{jjk}.
\]
Each summand on the first side is at most the corresponding summand on the
last side. Equality of the sums forces equality term by term. Positivity
of $N_{iik}$ implies that, for fixed nonunit $k$, all $t_{iik}$ have the
same sign $\epsilon_k$. The gauge $x_k\mapsto\epsilon_k x_k$ changes
$t_{iik}$ by the factor $\epsilon_i^2\epsilon_k=\epsilon_k$, proving the
claim. The unit coefficients are preserved.
\end{proof}

There are $25$ repeated-index unordered nonunit triples and ten triples
with all indices distinct. After the proved gauge, all $2^{10}=1024$
remaining sign assignments are checked exactly for associativity. Only
the all-positive assignment survives. Thus the signed algebra is a basis
sign gauge of $R(8)$. Of its six characters, the four codegree-five
characters have zero entries. Such characters cannot be dimensions of
simple objects. Only the FP character and its quadratic conjugate remain.
Every hypothetical categorification consequently has a pseudo-unitary
Galois conjugate.

The complete all-orders Gauss-family verifier excludes that possibility:
it forces a sum of $770$ roots of unity to equal
\[
-\frac{8^2+8\sqrt{1604}}2=-32-8\sqrt{401},
\]
whereas the quadratic Gauss coefficient bound at conductor $401$ requires
at least $400\cdot8=3200$ roots. Therefore $R(8)$ has no complex fusion
categorification.

The standalone executable \texttt{exclusion8200.py} reconstructs the
full family tensor, Casimir polynomial, integral-basis arithmetic, every
mod-two character, strengthened trace bound, every normalized sign table,
original dimension-character obstruction, and the complete existing Gauss
certificate. Its verifier recomputes the entire certificate and works
unchanged under optimized Python. The new proof does not use a signed
search timeout, a bounded root-of-unity order, or an assumption that
arbitrary fusion categories admit pivotal structures.

The categorical inputs are sphericalization and positivity from
Etingof--Nikshych--Ostrik, \emph{On fusion categories},
Ann. Math. 162 (2005), arXiv:math/0203060, and central induction from
Ostrik, \emph{Pivotal fusion categories of rank 3}, Theorems 2.5 and 2.13,
Mosc. Math. J. 15 (2015), arXiv:1309.4822. The new isolation, parity and
repeated-index gauge arguments are proved above.
